\chapter{Heller Cardiomyotomy}

\section*{Introduction \& Clinical Indications}
Heller cardiomyotomy, often referred to as Heller myotomy, is a surgical procedure designed to alleviate the symptoms of achalasia, a chronic esophageal motility disorder characterized by the failure of the lower esophageal sphincter (LES) to relax and the absence of effective peristalsis in the esophagus. This surgery involves the longitudinal division of the muscular layers of the distal esophagus and the LES, extending onto the gastric cardia, while preserving the underlying mucosa. The primary objective of this intervention is to relieve the functional obstruction at the gastroesophageal junction, which leads to significant dysphagia, regurgitation, and discomfort in affected patients. By surgically weakening the hypertonic sphincter, the myotomy facilitates the passage of food and liquids into the stomach, thereby improving the patient's quality of life. In contemporary practice, Heller myotomy is predominantly performed using minimally invasive laparoscopic techniques, often accompanied by a partial fundoplication to reduce the risk of postoperative gastroesophageal reflux disease (GERD).

\begin{itemize}
  \item Laparoscopic Heller Myotomy
  \item Esophagomyotomy
  \item Laparoscopic Esophagomyotomy with Fundoplication
  \item Modified Heller Procedure
  \item Extramucosal Myotomy
  \item Cardiomyotomy for Achalasia
\end{itemize}

The surgical principle behind Heller cardiomyotomy is based on the pathophysiology of achalasia, which is characterized by the degeneration of inhibitory neurons in the myenteric plexus, leading to a failure of the LES to relax and a lack of coordinated peristalsis in the esophagus. This results in a functional obstruction at the gastroesophageal junction, causing the esophagus to dilate and retain food, leading to dysphagia and regurgitation.

The rationale for the myotomy is to surgically disrupt the hypertonic circular muscle fibers of the LES, which are responsible for its sustained contraction. By performing an extramucosal myotomy, the surgeon longitudinally divides these muscle layers, extending from the distal esophagus across the gastroesophageal junction and onto the gastric cardia. This division effectively widens the esophageal outlet, allowing for the passage of food and liquids into the stomach. A critical aspect of modern Heller myotomy is the concomitant performance of a partial fundoplication, which serves to prevent postoperative gastroesophageal reflux disease (GERD) by providing a mechanical barrier while allowing the LES to function properly.

The evolution of surgical treatment for achalasia has a rich history, beginning with Ernst Heller's pioneering work in 1913, when he performed the first successful myotomy for achalasia in Germany. His original technique involved a "double Heller" approach, dividing both the anterior and posterior muscle layers of the distal esophagus. Although effective, this method was associated with higher morbidity due to extensive dissection.

Over the years, the technique evolved, leading to the adoption of a single anterior myotomy as the standard approach, which proved equally effective with fewer complications. The introduction of minimally invasive techniques in the early 1990s marked a significant advancement in the field. The first laparoscopic Heller myotomy was reported in 1991, demonstrating superior outcomes in terms of reduced postoperative pain, shorter hospital stays, and improved cosmetic results compared to open surgery. The recognition of the need for an anti-reflux procedure, such as a Dor or Toupet fundoplication, further refined the technique, solidifying laparoscopic Heller cardiomyotomy with partial fundoplication as the preferred surgical treatment for achalasia.

Heller cardiomyotomy is indicated for patients with confirmed achalasia who experience significant symptoms and for whom less invasive treatments have failed. The primary clinical indications include:

\begin{itemize}
  \item Persistent and severe dysphagia for solids and liquids, significantly impairing nutritional intake and quality of life.
  \item Frequent regurgitation of undigested food, leading to aspiration risk and recurrent respiratory infections.
  \item Significant weight loss due to the inability to consume adequate nutrition.
  \item Recurrent non-cardiac chest pain characteristic of esophageal spasm, unresponsive to medical therapy.
  \item Failure of endoscopic treatments, such as pneumatic dilation or botulinum toxin injection, to provide durable symptom relief.
  \item Primary treatment for younger patients desiring a long-term solution with high success rates.
  \item Objective diagnostic findings consistent with achalasia, such as a "bird's beak" appearance on barium swallow and characteristic manometric findings.
\end{itemize}

Heller cardiomyotomy is considered a primary treatment for achalasia, particularly for patients who are good surgical candidates and desire a durable solution. It is often recommended as a first-line intervention for younger patients or those with vigorous achalasia (Type III) due to its excellent long-term success rates. 

However, it also serves as a secondary or salvage procedure for patients who have experienced failure or recurrence of symptoms after less invasive treatments. While peroral endoscopic myotomy (POEM) has emerged as another primary treatment option, Heller myotomy remains a robust alternative, especially in cases where POEM is technically challenging or contraindicated.

\section*{Relevant Surgical Anatomy}
The surgical anatomy relevant to Heller cardiomyotomy encompasses the distal esophagus, gastroesophageal junction (GEJ), and the surrounding structures that contribute to the function of the lower esophageal sphincter (LES). The esophagus is a muscular conduit approximately 25 cm in length, composed of an outer longitudinal and an inner circular muscle layer. The distal esophagus traverses the esophageal hiatus of the diaphragm, where the crura of the diaphragm play a crucial role in maintaining the pressure of the LES.

Key anatomical structures include:

\begin{itemize}
  \item \textbf{Esophagus:} The distal 6-8 cm of the esophagus is targeted during the myotomy, focusing on the circular muscle layer.
  \item \textbf{Gastroesophageal Junction (GEJ):} The junction where the esophagus meets the stomach, identifiable by the transition from squamous to columnar epithelium.
  \item \textbf{Diaphragmatic Hiatus:} The opening in the diaphragm through which the esophagus passes, critical for maintaining LES function.
  \item \textbf{Phrenoesophageal Ligament:} A connective tissue structure anchoring the esophagus to the diaphragm, which must be divided for adequate exposure.
  \item \textbf{Vagus Nerves:} The anterior and posterior vagal trunks are essential for gastric motility and must be preserved during surgery to prevent complications such as gastroparesis.
  \item \textbf{Arterial Supply:} Blood supply to the distal esophagus is primarily from the left gastric artery and inferior phrenic arteries.
  \item \textbf{Venous Drainage:} The esophageal veins drain into the azygos and hemiazygos systems, with connections to the portal system.
  \item \textbf{Lymphatic Drainage:} Lymphatics drain into nodes along the left gastric artery and celiac axis.
\end{itemize}

A thorough understanding of these anatomical relationships is essential for the safe execution of the myotomy, particularly in preserving the vagus nerves and ensuring mucosal integrity.

\section*{Preoperative Workup \& Preparation}
Thorough preoperative preparation is essential for optimizing patient condition and minimizing risks. Patients are usually instructed to follow a clear liquid diet for 2-3 days prior to surgery to reduce aspiration risk.

Key preparatory steps include:

\begin{itemize}
  \item \textbf{Diagnostic Confirmation:} Confirm the diagnosis of achalasia through barium swallow and high-resolution esophageal manometry.
  \item \textbf{Endoscopy:} Perform upper endoscopy to rule out pseudoachalasia and assess for mucosal abnormalities.
  \item \textbf{Medical Optimization:} Conduct a comprehensive medical evaluation and obtain baseline laboratory tests.
  \item \textbf{Cardiac and Anesthesia Clearance:} Ensure cardiac fitness for general anesthesia and laparoscopic surgery.
  \item \textbf{Medication Review:} Review and adjust medications, particularly anticoagulants, to minimize bleeding risk.
  \item \textbf{Antibiotic Prophylaxis:} Administer intravenous antibiotics shortly before the incision.
  \item \textbf{Informed Consent:} Obtain detailed informed consent, ensuring the patient understands the procedure and its risks.
\end{itemize}

\textbf{Absolute Contraindications:}
\begin{itemize}
  \item Severe, uncorrectable cardiopulmonary disease.
  \item Active, uncontrolled esophageal perforation or severe infection.
  \item Uncorrected coagulopathy.
  \item Carcinoma of the gastroesophageal junction or distal esophagus.
\end{itemize}

\textbf{Relative Contraindications \& Precautions:}
\begin{itemize}
  \item Prior extensive upper abdominal surgery.
  \item Severe esophageal dilation with tortuosity.
  \item Advanced age or extreme frailty.
  \item Morbid obesity (BMI > 40 kg/m\textsuperscript{2}).
  \item Pregnancy.
  \item Active peptic ulcer disease or severe esophagitis.
  \item Significant hiatal hernia.
\end{itemize}

\section*{Surgical Technique \& Procedure}
The laparoscopic Heller cardiomyotomy is performed under general anesthesia, with the patient positioned supine in a reverse Trendelenburg position. The procedure begins with the creation of pneumoperitoneum, followed by the insertion of several small trocars into the upper abdomen.

Key steps of the procedure include:

\begin{itemize}
  \item \textbf{Access and Exposure:} Dissect the gastrohepatic ligament and phrenoesophageal membrane to mobilize the distal esophagus and gastroesophageal junction.
  \item \textbf{Identification of the LES:} Locate the lower esophageal sphincter by identifying the thickened circular muscle layer.
  \item \textbf{Myotomy Initiation:} Begin the myotomy on the anterior surface of the esophagus, approximately 2-3 cm proximal to the gastroesophageal junction.
  \item \textbf{Extramucosal Dissection:} Carefully divide the circular muscle fibers longitudinally until the submucosal layer is visible, extending distally across the gastroesophageal junction and onto the gastric cardia.
  \item \textbf{Mucosal Integrity Check:} Confirm the integrity of the esophageal mucosa using an intraoperative endoscope and bubble test.
  \item \textbf{Fundoplication:} Perform a partial fundoplication (Dor or Toupet) to prevent postoperative GERD.
  \item \textbf{Closure:} Inspect the surgical field for hemostasis, remove trocars, and close the abdominal incisions.
\end{itemize}

A nasogastric tube may be placed and typically removed within the first few hours postoperatively. Drains are generally not used unless there is a concern for mucosal injury or significant bleeding.

\section*{Complications \& Risks}
Heller cardiomyotomy carries potential risks, categorized into general surgical complications and procedure-specific complications.

\begin{itemize}
  \item \textbf{General Surgical Complications:}
    \begin{itemize}
      \item \textit{Bleeding:} Intraoperative or postoperative hemorrhage.
      \item \textit{Infection:} Surgical site or intra-abdominal infections.
      \item \textit{Anesthesia Complications:} Adverse reactions to general anesthesia.
      \item \textit{DVT and PE:} Formation of blood clots.
      \item \textit{Injury to Adjacent Organs:} Accidental injury during dissection.
    \end{itemize}
  \item \textbf{Procedure-Specific Risks:}
    \begin{itemize}
      \item \textit{Esophageal Perforation:} A serious complication requiring immediate repair.
      \item \textit{GERD:} New-onset or worsened reflux symptoms.
      \item \textit{Incomplete Myotomy:} Leading to persistent dysphagia.
      \item \textit{Esophageal Stricture:} Rarely occurring due to scarring.
      \item \textit{Vagal Nerve Injury:} Potentially leading to gastroparesis.
      \item \textit{Pneumothorax/Pneumomediastinum:} Usually self-limiting.
      \item \textit{Port Site Hernia:} Herniation through incision sites.
      \item \textit{Recurrence of Achalasia Symptoms:} Occurring in a small percentage of patients.
      \item \textit{Dysphagia to Solids/Liquids:} Due to fundoplication issues.
    \end{itemize}
\end{itemize}

\section*{Postoperative Care \& Recovery}
Postoperative recovery following a laparoscopic Heller cardiomyotomy typically involves close monitoring in the Post-Anesthesia Care Unit (PACU) for vital signs and pain management. Patients usually experience mild to moderate incisional pain, which is managed with analgesics. A nasogastric tube, if placed, is often removed within the first few hours postoperatively.

On the first postoperative day, a contrast swallow study is often performed to assess the integrity of the esophageal repair. If the study is clear, patients are started on a clear liquid diet, which is gradually advanced to soft foods and eventually a regular diet over the next 1-2 weeks. The typical hospital stay for an uncomplicated laparoscopic Heller myotomy is 1-2 days. 

Upon discharge, patients receive detailed instructions on diet progression, pain management, and activity restrictions. Strenuous activities and heavy lifting are usually restricted for 4-6 weeks. Full recovery and return to normal activities typically occur within 4-6 weeks.

During a Heller cardiomyotomy, the surgeon typically observes a dilated distal esophagus with a thickened LES region. The myotomy aims to confirm the diagnosis of achalasia and ensure no obstructing masses are present. The integrity of the esophageal mucosa is confirmed post-myotomy, ensuring no perforations have occurred. 

While no tissue is routinely resected for histopathological examination, the underlying pathology of achalasia involves the loss of inhibitory ganglion cells in the myenteric plexus. If biopsies are taken, pathologists would look for inflammation or fibrosis indicative of the disease process.

A successful postoperative course following Heller cardiomyotomy is characterized by significant improvement in dysphagia and overall quality of life. Patients typically experience mild incisional pain, which is well-controlled with analgesics. They advance their diet from clear liquids to soft foods over the first few days, with noticeable improvements in swallowing ability.

Most patients are discharged within 1-2 days, and over the first few weeks, they should experience a progressive resolution of dysphagia and an increase in appetite. While some may experience mild reflux symptoms, these are generally manageable with proton pump inhibitors. The small laparoscopic incisions heal well, and patients can usually resume normal activities within 4-6 weeks.

Postoperative care is crucial for monitoring recovery and managing potential complications. 

\begin{itemize}
  \item \textbf{Initial Clinic Visit:} Scheduled 2-4 weeks post-surgery to assess healing and symptom improvement.
  \item \textbf{Dietary Guidance:} Continued advice on eating habits and dietary progression.
  \item \textbf{Medication Management:} Proton pump inhibitors may be prescribed for reflux symptoms.
  \item \textbf{Activity Resumption:} Gradual return to normal activities, with restrictions on heavy lifting for 4-6 weeks.
  \item \textbf{Symptom Monitoring:} Patients should monitor for any recurrence of symptoms.
  \item \textbf{Further Diagnostics:} If symptoms recur, further investigations may be necessary.
\end{itemize}

Patients are advised to contact their surgeon if they experience severe abdominal pain, fever, persistent nausea or vomiting, or signs of wound infection.

\section*{Outcomes \& Prognosis}
Achalasia is a rare disorder, with an estimated incidence of 1-2 cases per 100,000 individuals annually and a prevalence of 10-15 cases per 100,000 population. It typically affects adults between 25 and 60 years of age, with no significant sex predilection. Heller cardiomyotomy is one of the most frequently performed treatments for achalasia, with thousands of procedures conducted annually. The overall mortality rate associated with laparoscopic Heller myotomy is low, while morbidity rates related to esophageal perforation and postoperative reflux are more significant but generally manageable. The procedure has a high success rate in relieving dysphagia, making it a cornerstone of achalasia management.

The long-term outcomes of laparoscopic Heller cardiomyotomy with partial fundoplication are excellent, with success rates ranging from 85\% to 95\%. Patients typically experience significant improvements in dysphagia and quality of life, with sustained relief often lasting over a decade. While recurrence of dysphagia occurs in approximately 5-10\% of patients, most can manage symptoms effectively with medical therapy. Prognostic factors influencing success include patient age, duration of symptoms, and the degree of esophageal dilation prior to surgery. Overall, Heller cardiomyotomy provides a durable and effective solution for the majority of patients with achalasia, significantly enhancing their quality of life.

\section*{References}
\begin{enumerate}
  \item MedlinePlus. Heller Cardiomyotomy. U.S. National Library of Medicine; 2024.
  \item Brunicardi FC, Andersen DK, Billiar TR, et al., editors. Schwartz's Principles of Surgery. 12th ed. McGraw-Hill Education; 2024.
  \item Kahrilas PJ, Bredenoord AJ, Fox M, et al. The Chicago Classification of Esophageal Motility Disorders, v3.0. Neurogastroenterol Motil. 2015;27(2):160-174.
  \item American College of Surgeons. Surgical Education and Self-Assessment Program (SESAP\textsuperscript{\textregistered}) 18. American College of Surgeons; 2024.
\end{enumerate}

\clearpage
